\documentclass[12pt]{article}
\usepackage{amsmath,amssymb,amsfonts}
\usepackage[margin=1in]{geometry}

\begin{document}

\textbf{Chapter 8, Problem 7.}

\textbf{(a)} Suppose that we have an attractive $\delta$-function potential in one dimension, $V(x) = -P\delta(x)$. Using the result of Problem 7.10, show that this potential will support one and only one bound state, no matter how large $P$ may be. Determine the binding energy and wave function of this state.

\textbf{(b)} Suppose that particles of positive energy $E$ are incident on a $\delta$-function potential, $V(x) = P\delta(x)$, in one dimension. What fraction of the particles is transmitted through the potential barrier and what fraction is reflected?

\bigskip
\textbf{Solution.}

\textbf{Part (a):}
For a $\delta$-function potential $V(x) = -P\delta(x)$ with $P>0$ (attractive), the time-independent Schr\"odinger equation in one dimension is:
\[
-\frac{\hbar^2}{2m}\frac{d^2\psi}{dx^2} - P\delta(x)\psi = E\psi.
\]

For a bound state, $E = -E_B < 0$. Define $\kappa = \sqrt{2mE_B}/\hbar$. Away from $x=0$:
\[
\psi''= \kappa^2\psi,
\]
giving $\psi(x) = Ae^{-\kappa|x|}$ (normalizable solution).

Integrating the Schr\"odinger equation across $x=0$:
\[
-\frac{\hbar^2}{2m}[\psi'(0^+)-\psi'(0^-)] - P\psi(0) = 0.
\]
With $\psi'(0^+) = -\kappa A$ and $\psi'(0^-) = \kappa A$:
\[
-\frac{\hbar^2}{2m}(-2\kappa A) = PA \implies \kappa = \frac{mP}{\hbar^2}.
\]

This gives exactly \emph{one} bound state with energy:
\[
\boxed{E_B = \frac{\hbar^2\kappa^2}{2m} = \frac{mP^2}{2\hbar^2},}
\]
and normalized wave function:
\[
\boxed{\psi(x) = \sqrt{\kappa}\,e^{-\kappa|x|} = \sqrt{\frac{mP}{\hbar^2}}\,e^{-mP|x|/\hbar^2}.}
\]

Since the matching condition gives a unique $\kappa$, there is exactly one bound state regardless of the value of $P$.

\bigskip
\textbf{Part (b):}
For scattering with $V(x) = P\delta(x)$ and incident energy $E>0$, let $k = \sqrt{2mE}/\hbar$. The wave function is:
\[
\psi(x) = \begin{cases} e^{ikx} + Re^{-ikx}, & x<0, \\ Te^{ikx}, & x>0. \end{cases}
\]

Continuity at $x=0$: $1+R = T$.

Discontinuity of derivative (integrating across $x=0$):
\[
-\frac{\hbar^2}{2m}[ikT - ik(1-R)] + PT = 0,
\]
\[
-\frac{\hbar^2}{2m}ik(T-1+R) = -PT.
\]
Using $T = 1+R$:
\[
-\frac{\hbar^2}{2m}ik(2R) = -P(1+R) \implies -i\hbar^2 kR/m = -P(1+R).
\]
Solving:
\[
R = \frac{-P}{P + i\hbar^2 k/m} = \frac{-mP}{mP+i\hbar^2 k}.
\]
\[
T = 1+R = \frac{i\hbar^2 k}{mP+i\hbar^2 k}.
\]

The transmission and reflection coefficients:
\[
\boxed{|T|^2 = \frac{\hbar^4 k^2}{m^2P^2+\hbar^4 k^2} = \frac{1}{1+m^2P^2/(\hbar^4 k^2)} = \frac{1}{1+mP^2/(2\hbar^2 E)},}
\]
\[
\boxed{|R|^2 = \frac{m^2P^2}{m^2P^2+\hbar^4 k^2} = \frac{1}{1+2\hbar^2 E/(mP^2)}.}
\]

One can verify $|R|^2+|T|^2=1$. $\blacksquare$

\end{document}
